
\documentclass[os,manuscript]{copernicus}
\usepackage{booktabs}
\usepackage{amsmath}
\begin{document}
\title{Copernicus Family: Official Copernicus Submission Preview}
\Author[1][verified@example.org]{A.}{Placeholder}
\Author[1]{B.}{Example}
\Author[2]{C.}{Template}
\affil[1]{Department of Verified Templates, Placeholder Institute, City, Country}
\affil[2]{Journal Systems Lab, Example University, City, Country}
\runningtitle{Copernicus Family: Official Copernicus Submission Preview}
\runningauthor{Placeholder et al.}
\begin{abstract}
This submission-style preview is rendered with the imported official Copernicus package. The page is intentionally structured as a realistic geoscience manuscript so front matter, section rhythm, figure-table balance, and declaration blocks can be checked against the family baseline.
\end{abstract}
\section{Introduction}
The introduction is written as compact technical prose so the preview reads more like a practical Copernicus submission than a generic display page.
\section{Related Work}
The related-work section is explicit so the preview reveals how literature framing, methods, and experiments separate within the family baseline.
\section{Methodology}
The methodology section introduces representative mathematics and a figure placeholder so the local Copernicus package can be inspected under more realistic technical content.
\begin{equation}
\mathcal{L} = \sum_{i=1}^{M} w_i \left\| y_i - \hat{y}_i \right\|_2^2.
\end{equation}
\section{Experiments}
The experiments section provides enough structure to expose table treatment and section rhythm inside the Copernicus family baseline.
\section{Conclusion and Discussion}
This page verifies the imported official Copernicus package path while preserving the six-part manuscript skeleton requested for the library.
\authorcontribution{All contributions are placeholders for template verification.}
\competinginterests{The authors declare that they have no conflict of interest.}
\begin{thebibliography}{9}
\bibitem[Placeholder and Example(2026)]{ref1} Placeholder, A. and Example, B.: Template-aware Copernicus manuscript preparation, Placeholder Journal, 1, 1--2, 2026.
\end{thebibliography}
\end{document}
